\documentclass[12pt,a4paper]{article}
\usepackage{ugmath}
\usepackage{float}
\usepackage{placeins}
\usepackage{array}
\usepackage[labelfont=bf,labelsep=space]{caption}
\newcommand{\paren}[1]{\left( #1 \right)}
\newcommand{\alumno}{Ricardo León Martínez}
\newcommand{\materia}{Elementos de Probabilidad y Estadistica}
\newcommand{\profesor}{Claudia Reynoso Alcántara}
\newcommand{\tarea}{Tarea 10}
\newcommand{\fecha}{22/5/2026}

\begin{document}

    \begin{center}
        {\large\textbf{UNIVERSIDAD DE GUANAJUATO}}\\[0.3cm]
        {\normalsize\textbf{DIVISIÓN DE CIENCIAS NATURALES Y EXACTAS}}\\
        {\normalsize\textbf{CAMPUS GUANAJUATO}}\\[1cm]

        {\Large\textbf{\tarea\ (\materia)}}\\[1cm]
    \end{center}

    \noindent
    \textbf{Nombre:} \alumno 
    \hfill 
    \textbf{Fecha:} \fecha 
    \hfill 
    \textbf{Calificación:} \rule{3cm}{0.4pt}

    \vspace{0.3cm}

    \section*{Parte I. Valor esperado}

    \begin{mdframed}[style=mdpinkbox, frametitle={Ejercicio 1}]
        Sea $X$ una variable aleatoria discreta con soporte contenido en una sucesión $\{a_i\}_{i=1}^{\infty}$
        con $a_i = a_j$ si y sólo si $i = j$, es decir, se cumple que
        \begin{equation*}
            1 = \sum_{i=1}^{\infty} P(X = a_i).
        \end{equation*}
        Fijemos un entero $k \geq 1$ y consideremos las variables aleatorias $X^k$, $(X+1)^k$ y $(X-1)^k$.
        Escriba expresiones para
        \begin{equation*}
            \mathbb{E}[X^k], \quad \mathbb{E}[(X+1)^k] \quad \text{y} \quad \mathbb{E}[(X-1)^k]
        \end{equation*}
        como series. Puede suponer que cada uno de los valores esperados anteriores existen.
    \end{mdframed}

    \begin{mdframed}[style=mdpinkbox, frametitle={Ejercicio 2}]
        Sea $X$ una variable aleatoria discreta con soporte contenido en los naturales, es decir, se cumple que
        \begin{equation*}
            1 = \sum_{n=1}^{\infty} P(X = n).
        \end{equation*}
        Demuestre que
        \begin{equation*}
            \mathbb{E}[X] = \sum_{k=1}^{\infty} P(X \geq k).
        \end{equation*}
    \end{mdframed}

    \section*{Parte II. Distribución binomial}

    Recordemos que una variable aleatoria $X$ tiene distribución binomial de parámetros
    $(n, p) \in \mathbb{N} \times (0,1)$ si se cumple que
    \begin{equation*}
        P(X = i) = \binom{n}{i} p^i (1-p)^{n-i} \quad \text{para } i = 0, 1, \ldots, n.
    \end{equation*}

    \begin{mdframed}[style=mdpinkbox, frametitle={Ejercicio 3}]
        Sean $X$ e $Y$ variables aleatorias binomiales de parámetros $(n, p)$ y $(n-1, p)$,
        respectivamente. Demuestre que para todo entero $k \geq 1$ se cumple que
        \begin{equation*}
            \mathbb{E}[X^k] = np \cdot \mathbb{E}[(Y+1)^{k-1}].
        \end{equation*}
        Encuentre $E[X]$ y $E[X^2]$ y concluya que $\mathrm{Var}(X) = np(1-p)$.

        \medskip
        \noindent\textit{Sugerencia.} Escriba $\mathbb{E}[X^k]$ como una serie y considere que
        $i\binom{n}{i} = n\binom{n-1}{i-1}$.
    \end{mdframed}

    \section*{Parte III. Distribución geométrica}

    Una variable aleatoria $X$ tiene distribución geométrica de parámetro $p \in (0,1)$ si se cumple que
    \begin{equation*}
        P(X = i) = p(1-p)^{i-1} \quad \text{para } i = 0, 1, 2, \ldots
    \end{equation*}

    \begin{mdframed}[style=mdpinkbox, frametitle={Ejercicio 4}]
        Sea $X$ una variable aleatoria geométrica de parámetro $p \in (0,1)$. Tomemos $q = 1 - p$.
        Demuestre que
        \begin{align*}
            \mathbb{E}[X]   &= q\,\mathbb{E}[X] + 1, \\
            \mathbb{E}[X^2] &= q\,\mathbb{E}[X^2] + 2q\,\mathbb{E}[X] + 1.
        \end{align*}
        Encuentre $E[X]$ y $E[X^2]$ y concluya que $\mathrm{Var}(X) = \dfrac{1-p}{p^2}$.

        \medskip
        \noindent\textit{Sugerencias.} Escriba $\mathbb{E}[X]$ como una serie y considere que
        $i q^{i-1} p = (i-1)q^{i-1}p + q^{i-1}p$. De manera similar, escriba $\mathbb{E}[X^2]$
        como una serie y considere que
        \begin{equation*}
            i^2 q^{i-1} p = (i-1)^2 q^{i-1}p + 2(i-1)q^{i-1}p + q^{i-1}p.
        \end{equation*}
    \end{mdframed}

    \section*{Parte IV. Distribución binomial negativa}

    Una variable aleatoria $X$ tiene distribución binomial negativa de parámetros $(r, p)$ si se cumple que
    \begin{equation*}
        P(X = i) = \binom{i-1}{r-1} p^r (1-p)^{i-r} \quad \text{para } i = r, r+1, r+2, \ldots
    \end{equation*}

    \begin{mdframed}[style=mdpinkbox, frametitle={Ejercicio 5}]
        Sean $X$ e $Y$ variables aleatorias binomiales negativas de parámetros $(r, p)$ y $(r+1, p)$,
        respectivamente. Demuestre que para todo entero $k \geq 1$ se cumple que
        \begin{equation*}
            \mathbb{E}[X^k] = \frac{r}{p} \cdot \mathbb{E}[(Y-1)^{k-1}].
        \end{equation*}
        Encuentre $E[X]$ y $E[X^2]$ y concluya que $\mathrm{Var}(X) = r \cdot \dfrac{1-p}{p^2}$.

        \medskip
        \noindent\textit{Sugerencia.} Escriba $\mathbb{E}[X^k]$ como una serie y considere que
        \begin{equation*}
            i^k \binom{i-1}{r-1} p^r = \frac{r}{p}\, i^{k-1} \binom{i}{r} p^{r+1}.
        \end{equation*}
    \end{mdframed}

    \section*{Parte V. Distribución hipergeométrica}

    Recordemos que una variable aleatoria $X$ tiene distribución hipergeométrica de parámetros $(b, n, m)$
    con $m \leq b + n$ si se cumple que
    \begin{equation*}
        P(X = i) = \frac{\dbinom{b}{i}\dbinom{n}{m-i}}{\dbinom{b+n}{m}} \quad \text{para } i = 0, 1, \ldots, m,
    \end{equation*}
    donde tomamos $\binom{b}{i} = 0$ si $i > b$ y $\binom{n}{m-i} = 0$ si $m - i > n$.

    \begin{mdframed}[style=mdpinkbox, frametitle={Ejercicio 6}]
        Sean $X$ e $Y$ variables aleatorias hipergeométricas de parámetros $(b, n, m)$ y
        $(b-1, n, m-1)$, respectivamente. Demuestre que para todo entero $k \geq 1$ se cumple que
        \begin{equation*}
            \mathbb{E}[X^k] = m\,\frac{b}{b+n} \cdot \mathbb{E}[(Y+1)^{k-1}].
        \end{equation*}
        Encuentre $E[X]$ y $E[X^2]$. Tome $p = \dfrac{b}{b+n}$ y concluya que
        \begin{equation*}
            \mathrm{Var}(X) = mp(p-1)\!\left(1 - \frac{m-1}{b+n-1}\right).
        \end{equation*}

        \medskip
        \noindent\textit{Sugerencia.} Escriba $\mathbb{E}[X^k]$ como una serie y considere que
        $i^k \binom{b}{i} = i^{k-1} b \binom{b-1}{i-1}$ y que $m\binom{b+n}{m} = (b+n)\binom{b+n-1}{m-1}$.
    \end{mdframed}

    \section*{Parte VI. Distribución Poisson}

    Una variable aleatoria $X$ tiene distribución Poisson de parámetro $\lambda > 0$ si se cumple que
    \begin{equation*}
        P(X = i) = \frac{\lambda^i}{i!} e^{-\lambda} \quad \text{para } i = 0, 1, 2, \ldots
    \end{equation*}
    El objetivo del siguiente ejercicio es ver que una variable aleatoria Poisson puede considerarse como
    límite de variables aleatorias binomiales.

    \begin{mdframed}[style=mdpinkbox, frametitle={Ejercicio 7}]
        Sea $X$ una variable aleatoria Poisson de parámetro $\lambda > 0$. Supongamos que $X_n$ es una
        variable aleatoria binomial de parámetros $(n, p_n)$ con $p_n = \lambda/n$ para $n = 1, 2, \ldots$
        Demuestre que para todo entero $i \geq 0$ se cumple que
        \begin{equation*}
            P(X = i) = \lim_{n \to \infty} P(X_n = i).
        \end{equation*}

        \medskip
        \noindent\textit{Sugerencia.} Puede dar por hecho tanto que
        $e^t = \sum_{i=0}^{\infty} t^i / i!$ como que $e^t = \lim_{n\to\infty}(1 + t/n)^n$
        para todo $t \in \mathbb{R}$.
    \end{mdframed}

    \begin{mdframed}[style=mdpinkbox, frametitle={Ejercicio 8}]
        Sea $X$ una variable aleatoria Poisson de parámetro $\lambda > 0$. Encuentre $E[X]$ y $E[X^2]$
        y concluya que $\mathrm{Var}(X) = \lambda$.

        \medskip
        \noindent\textit{Sugerencia.} Note que
        \begin{equation*}
            i\,\frac{\lambda^i}{i!} = \lambda\,\frac{\lambda^{i-1}}{(i-1)!}
            \qquad \text{y} \qquad
            i^2\,\frac{\lambda^i}{i!} = \lambda^2\,\frac{\lambda^{i-2}}{(i-2)!} + \lambda\,\frac{\lambda^{i-1}}{(i-1)!}.
        \end{equation*}
    \end{mdframed}
\end{document}